\documentclass[11pt]{article}
\usepackage[margin=1in]{geometry}
\usepackage{amsmath,amssymb,amsthm,mathtools,bm}
\usepackage{microtype}
\usepackage[T1]{fontenc}
\usepackage{lmodern}
\usepackage{booktabs,enumitem}
\usepackage[hidelinks]{hyperref}
\setlist{nosep}
\newtheorem{theorem}{Theorem}[section]
\newtheorem{corollary}[theorem]{Corollary}
\newtheorem{proposition}[theorem]{Proposition}
\theoremstyle{definition}
\newtheorem{definition}[theorem]{Definition}
\theoremstyle{remark}
\newtheorem{remark}[theorem]{Remark}

\newcommand{\dd}{\mathrm{d}}

\title{\textbf{Synopsis: Hysterical Conformal Cyclic Cosmology}\\[0.35em]
\large Infrared electromagnetic criticality and an active big-bang crossover}
\author{Andrew Bruce Boyd}
\date{September 2026}

\begin{document}
\maketitle
\noindent\textit{Computational note.}
Proofs, exploratory experiments, and paper writing were assisted by generative AI.
Mathematical theorems stated in this paper are formally verified in Lean~4 in the
companion specifications, as faithful ports of the TeX arguments at the abstractions
recorded there, except results explicitly tagged as \emph{literature hypotheses}
(standard theorems cited from the literature and not re-proved here) or
\emph{physical inputs} (modeling assumptions outside mathematics).

\begin{abstract}
This is a short guide to the Hysterical Conformal Cyclic Cosmology (HCCC) paper.
Each section says \emph{why} the claim matters, then states the result.
Proofs and the full status table live in the full paper.
\end{abstract}

\section{Why an electromagnetic late-universe continuation}
The late-universe gravitational response forces dilution criticality and a vacuum-like
attractor.
Charge conservation still forbids creating homogeneous net charge from neutrality, but
finite-wavenumber electromagnetic response remains allowed.
The question here is whether that infrared electromagnetic channel can drive an
\emph{active} conformal crossover---a big-bang-like event for the next aeon---in the
geometric spirit of Penrose's conformal cyclic cosmology, but with a dynamical phase
transition rather than a passive identification.

\section{Linear infrared instability}
For a helical electromagnetic--response mode with physical wavenumber $q$ in de~Sitter
freezing, the characteristic polynomial is
\[
p_q(\lambda)=(\lambda+\gamma)\bigl((\lambda+2H)^2+q^2\bigr)-A(\lambda+2H).
\]

\begin{theorem}[Zero-mode stationary threshold]
$p_0(0)=0$ if and only if $A=2H\gamma$.
\end{theorem}

\begin{theorem}[Infrared-first stationary criticality]
For positive parameters, the stationary gain decreases in $q^2$; the $q=0$ mode reaches
the stationary threshold first.
\end{theorem}

\begin{theorem}[Infrared unstable band]
$p_q(0)<0$ if and only if $q^2<2HA/\gamma-4H^2$.
When $A>2H\gamma$ this is a nonempty infrared band, and $p_q(0)<0$ guarantees a positive
real eigenvalue.
\end{theorem}

Near threshold the growth rate is maximal at $q=0$: an infrared onset, not an ultraviolet
catastrophe.

\section{Photon pump (modelling)}
A scalar response with Maxwell kinetic function $I(h)$ yields the mode equation
$u_k''+(k^2-I''/I)u_k=0$.
A non-adiabatic $h$ therefore pumps finite-$k$ photons.
Efficiency, backreaction, and saturation are open dynamical problems, not theorems.

\section{Reciprocal conformal matching}
Across a crossover with reciprocal conformal factors $\Omega_+\Omega_-=C_0$, four-dimensional
Maxwell theory transports fixed $F_{\mu\nu}$ while densities scale as $\rho=\Omega^{-4}\hat\rho$.

\begin{theorem}[Cold outgoing / hot incoming]
Finite conformal radiation density appears cold on the outgoing side and hot on the
incoming side; with unit reciprocity, $\rho_+/\rho_-=a_-^{8}$.
\end{theorem}

\begin{theorem}[Reciprocal de~Sitter $\to$ radiation]
Unit reciprocity maps outgoing $a_-\sim 1/(H_-|\eta|)$ to incoming $a_+\sim H_-\eta$, hence
$a_+\propto t_+^{1/2}$.
\end{theorem}

The hot incoming branch is the big-bang-like content of the next aeon.

\section{Active versus passive CCC}
Penrose CCC and HCCC share the conformal aeon-bridge geometry.
CCC is passive after matter becomes irrelevant; HCCC is active: an infrared Hysterical
electromagnetic phase transition produces the conformal radiation seed.
Whether our observed Big Bang was such a crossover is a physical hypothesis, not a theorem.

\section{What is not claimed}
No derivation of the saturated amplitude $\hat\rho_{\mathrm{EM}}(\Sigma)$; no complete
Einstein--Maxwell--response junction; no entropy/thermalisation proof; no observational CCC
ring or Hawking-point claim; no claim that Nature realises HCCC.
\end{document}
